\documentclass[11pt,a4paper,modern]{FFExercises}
\usepackage[english,greek]{babel}
\usepackage[utf8]{inputenc}
\usepackage{nimbusserif}
\usepackage[T1]{fontenc}
\usepackage{amsmath}
\let\myBbbk\Bbbk
\let\Bbbk\relax
\usepackage[amsbb,subscriptcorrection,zswash,mtpcal,mtphrb,mtpfrak]{mtpro2}
\usepackage{graphicx,multicol,multirow,enumitem,tabularx,mathimatika,gensymb,venndiagram,hhline,longtable,tkz-euclide,fontawesome5,eurosym,tcolorbox,tabularray,tikzpagenodes,relsize,siunitx,eurosym}
\definecolor{xrwma}{HTML}{aa1212}
\usetikzlibrary{calc}
\usetikzlibrary{positioning}
\tcbuselibrary{skins,theorems,breakable}
\renewcommand{\textstigma}{\textsigma\texttau}
\renewcommand{\textdexiakeraia}{}
\usepackage{svg}
\ekthetesdeiktes

\begin{document}

\titlos{Μαθηματικά Προσανατολισμού}{Γ΄ Λυκείου}{Ακρότατα Συνάρτησης (Κριτήριο Παραγώγου)}

%=============================================
% ΟΜΑΔΑ Α: ΥΠΟΛΟΓΙΣΤΙΚΕΣ - ΠΟΛΥΩΝΥΜΙΚΕΣ
%=============================================
\paragraph{Ομάδα Α: Υπολογιστικές Ασκήσεις - Πολυωνυμικές Συναρτήσεις}

\askhsh
Δίνεται η συνάρτηση $f(x) = x^3 - 3x^2 + 5$.
\begin{alist}
    \item Να βρείτε το πεδίο ορισμού της $f$ και να αποδείξετε ότι είναι συνεχής.
    \item Να υπολογίσετε την $f'(x)$.
    \item Να βρείτε τα κρίσιμα σημεία και να μελετήσετε τη μονοτονία.
    \item Να βρείτε τα τοπικά ακρότατα της $f$.
    \item Να εξετάσετε αν η εξίσωση $f(x) = 3$ έχει λύση.
\end{alist}

\askhsh
Δίνεται η συνάρτηση $f(x) = -2x^3 + 9x^2 - 12x + 4$.
\begin{alist}
    \item Να υπολογίσετε την $f'(x)$.
    \item Να λύσετε την εξίσωση $f'(x) = 0$.
    \item Να μελετήσετε τη μονοτονία της $f$.
    \item Να βρείτε τα τοπικά ακρότατα και τις τιμές τους.
    \item Να σχεδιάσετε τον πίνακα προσήμου - μονοτονίας.
\end{alist}

\askhsh
Δίνεται η συνάρτηση $f(x) = x^4 - 4x^3 + 10$.
\begin{alist}
    \item Να υπολογίσετε την $f'(x)$.
    \item Να βρείτε τα σημεία στα οποία μηδενίζεται η $f'$.
    \item Να εξετάσετε αν στα σημεία αυτά παρουσιάζονται ακρότατα.
    \item Να βρείτε τα διαστήματα μονοτονίας της $f$.
    \item Να υπολογίσετε τα $\lim_{x \to \pm\infty} f(x)$.
\end{alist}

\askhsh
Δίνεται η συνάρτηση $f(x) = 3x^4 - 8x^3 + 6x^2 + 1$.
\begin{alist}
    \item Να υπολογίσετε την $f'(x)$ και να την παραγοντοποιήσετε.
    \item Να βρείτε τα κρίσιμα σημεία της $f$.
    \item Να εξετάσετε αν υπάρχουν ακρότατα στα κρίσιμα σημεία.
    \item Να βρείτε την ελάχιστη τιμή της $f$.
\end{alist}

\askhsh
Δίνεται η συνάρτηση $f(x) = x^3 - 6x^2 + 9x + 2$.
\begin{alist}
    \item Να βρείτε την $f'(x)$.
    \item Να λύσετε την $f'(x) = 0$.
    \item Να εξετάσετε τι συμβαίνει στο σημείο μηδενισμού της παραγώγου.
    \item Να προσδιορίσετε αν η $f$ είναι γνησίως μονότονη.
\end{alist}

\askhsh
Δίνεται η συνάρτηση $f(x) = x^5 - 5x^3$.
\begin{alist}
    \item Να δείξετε ότι η $f$ είναι περιττή συνάρτηση.
    \item Να υπολογίσετε την $f'(x)$ και να βρείτε τα κρίσιμα σημεία.
    \item Να μελετήσετε τη μονοτονία και τα τοπικά ακρότατα.
    \item Να βρείτε τα $\lim_{x \to \pm\infty} f(x)$ και να εξηγήσετε γιατί δεν υπάρχουν ολικά ακρότατα.
    \item Να βρείτε το πλήθος των ριζών της εξίσωσης $f(x) = 0$.
\end{alist}

%=============================================
% ΟΜΑΔΑ Β: ΥΠΟΛΟΓΙΣΤΙΚΕΣ - ΡΗΤΕΣ
%=============================================
\paragraph{Ομάδα Β: Υπολογιστικές Ασκήσεις - Ρητές Συναρτήσεις}

\askhsh
Δίνεται η συνάρτηση $f(x) = x + \frac{1}{x}$.
\begin{alist}
    \item Να βρείτε το πεδίο ορισμού της $f$.
    \item Να υπολογίσετε την $f'(x)$.
    \item Να βρείτε τα κρίσιμα σημεία.
    \item Να μελετήσετε τη μονοτονία σε κάθε διάστημα του πεδίου ορισμού.
    \item Να βρείτε τα τοπικά ακρότατα.
\end{alist}

\askhsh
Δίνεται η συνάρτηση $f(x) = \frac{x}{x^2+1}$.
\begin{alist}
    \item Να βρείτε το πεδίο ορισμού.
    \item Να υπολογίσετε την $f'(x)$ χρησιμοποιώντας τον κανόνα του πηλίκου.
    \item Να βρείτε τα κρίσιμα σημεία.
    \item Να μελετήσετε τη μονοτονία.
    \item Να βρείτε τα ακρότατα και το σύνολο τιμών της $f$.
\end{alist}

\askhsh
Δίνεται η συνάρτηση $f(x) = \frac{x^2 - 2x + 1}{x^2 + 1}$.
\begin{alist}
    \item Να βρείτε το πεδίο ορισμού.
    \item Να υπολογίσετε την $f'(x)$.
    \item Να βρείτε τα κρίσιμα σημεία.
    \item Να μελετήσετε τη μονοτονία και τα ακρότατα.
    \item Να βρείτε το σύνολο τιμών της $f$.
\end{alist}

\askhsh
Δίνεται η συνάρτηση $f(x) = \frac{x^2}{x+1}$.
\begin{alist}
    \item Να βρείτε το πεδίο ορισμού.
    \item Να υπολογίσετε την $f'(x)$.
    \item Να βρείτε τα κρίσιμα σημεία.
    \item Να μελετήσετε τη μονοτονία.
    \item Να βρείτε τα τοπικά ακρότατα.
\end{alist}

\askhsh
Δίνεται η συνάρτηση $f(x) = \frac{x^2 - 4}{x^2 + 4}$.
\begin{alist}
    \item Να βρείτε το πεδίο ορισμού.
    \item Να υπολογίσετε την $f'(x)$.
    \item Να βρείτε τα κρίσιμα σημεία.
    \item Να μελετήσετε τη μονοτονία και τα ακρότατα.
\end{alist}

\askhsh
Δίνεται η συνάρτηση $f(x) = \frac{2x}{x^2 - 1}$.
\begin{alist}
    \item Να βρείτε το πεδίο ορισμού.
    \item Να υπολογίσετε την $f'(x)$.
    \item Να εξετάσετε αν υπάρχουν κρίσιμα σημεία.
    \item Να μελετήσετε τη μονοτονία σε κάθε διάστημα του πεδίου ορισμού.
\end{alist}

%=============================================
% ΟΜΑΔΑ Γ: ΥΠΟΛΟΓΙΣΤΙΚΕΣ - ΕΚΘΕΤΙΚΕΣ/ΛΟΓΑΡΙΘΜΙΚΕΣ
%=============================================
\paragraph{Ομάδα Γ: Υπολογιστικές Ασκήσεις - Εκθετικές και Λογαριθμικές}

\askhsh
Δίνεται η συνάρτηση $f(x) = x \cdot e^x$.
\begin{alist}
    \item Να βρείτε το πεδίο ορισμού.
    \item Να υπολογίσετε την $f'(x)$.
    \item Να βρείτε τα κρίσιμα σημεία.
    \item Να μελετήσετε τη μονοτονία.
    \item Να βρείτε τα ακρότατα.
\end{alist}

\askhsh
Δίνεται η συνάρτηση $f(x) = x \cdot e^{-x}$.
\begin{alist}
    \item Να υπολογίσετε την $f'(x)$.
    \item Να βρείτε τα κρίσιμα σημεία.
    \item Να μελετήσετε τη μονοτονία.
    \item Να βρείτε τα ακρότατα.
    \item Να υπολογίσετε τα $\lim_{x \to \pm\infty} f(x)$.
\end{alist}

\askhsh
Δίνεται η συνάρτηση $f(x) = \frac{\ln x}{x}, \quad x > 0$.
\begin{alist}
    \item Να υπολογίσετε την $f'(x)$.
    \item Να βρείτε τα κρίσιμα σημεία.
    \item Να μελετήσετε τη μονοτονία.
    \item Να βρείτε τα ακρότατα.
    \item Να υπολογίσετε τα $\lim_{x \to 0^+} f(x)$ και $\lim_{x \to +\infty} f(x)$.
\end{alist}

\askhsh
Δίνεται η συνάρτηση $f(x) = x^2 \ln x, \quad x > 0$.
\begin{alist}
    \item Να υπολογίσετε την $f'(x)$.
    \item Να βρείτε τα κρίσιμα σημεία.
    \item Να μελετήσετε τη μονοτονία.
    \item Να βρείτε τα ακρότατα.
\end{alist}

\askhsh
Δίνεται η συνάρτηση $f(x) = e^x - x - 1$.
\begin{alist}
    \item Να υπολογίσετε την $f'(x)$.
    \item Να βρείτε τα κρίσιμα σημεία.
    \item Να μελετήσετε τη μονοτονία και τα ακρότατα.
    \item Να δείξετε ότι $f(x) \ge 0$ για κάθε $x \in \mathbb{R}$.
\end{alist}

\askhsh
Δίνεται η συνάρτηση $f(x) = \ln x - x + 1, \quad x > 0$.
\begin{alist}
    \item Να υπολογίσετε την $f'(x)$.
    \item Να βρείτε τα κρίσιμα σημεία.
    \item Να μελετήσετε τη μονοτονία και τα ακρότατα.
    \item Να δείξετε ότι $\ln x \le x - 1$ για κάθε $x > 0$.
\end{alist}

%=============================================
% ΟΜΑΔΑ Δ: ΠΑΡΑΜΕΤΡΙΚΕΣ ΑΣΚΗΣΕΙΣ
%=============================================
\paragraph{Ομάδα Δ: Παραμετρικές Ασκήσεις}

\askhsh
Δίνεται η συνάρτηση $f(x) = x^2 - 2\alpha x + 3$.
\begin{alist}
    \item Να βρείτε την $f'(x)$.
    \item Να βρείτε το κρίσιμο σημείο συναρτήσει του $\alpha$.
    \item Να βρείτε το $\alpha$ ώστε η $f$ να παρουσιάζει ελάχιστο στο $x_0 = 2$.
    \item Για την τιμή αυτή του $\alpha$, να υπολογίσετε την ελάχιστη τιμή.
\end{alist}

\askhsh
Δίνεται η συνάρτηση $f(x) = \alpha x^3 + \beta x^2 - 6x + 1$.
\begin{alist}
    \item Να βρείτε την $f'(x)$.
    \item Να βρείτε τα $\alpha, \beta$ ώστε η $f$ να παρουσιάζει ακρότατα στα $x_1 = 1$ και $x_2 = -2$.
    \item Για τις τιμές αυτές, να προσδιορίσετε το είδος των ακροτάτων.
    \item Να υπολογίσετε τις τιμές $f(1)$ και $f(-2)$.
\end{alist}

\askhsh
Δίνεται η συνάρτηση $f(x) = x^3 - 3\alpha x^2 + 4$, όπου $\alpha \neq 0$.
\begin{alist}
    \item Να βρείτε την $f'(x)$.
    \item Να βρείτε τα κρίσιμα σημεία συναρτήσει του $\alpha$.
    \item Να βρείτε το $\alpha$ ώστε η $f$ να παρουσιάζει ακρότατο στο $x_0 = 2$.
    \item Για αυτή την τιμή, να μελετήσετε πλήρως τη μονοτονία.
\end{alist}

\askhsh
Δίνεται η συνάρτηση $f(x) = e^x - \alpha x$.
\begin{alist}
    \item Να βρείτε την $f'(x)$.
    \item Να βρείτε το κρίσιμο σημείο συναρτήσει του $\alpha > 0$.
    \item Να βρείτε το $\alpha$ ώστε η $f$ να παρουσιάζει ελάχιστο στο $x_0 = 0$.
    \item Για αυτή την τιμή του $\alpha$, να βρείτε την ελάχιστη τιμή.
\end{alist}

\askhsh
Δίνεται η συνάρτηση $f(x) = \ln x - \lambda x, \quad x > 0$.
\begin{alist}
    \item Να βρείτε την $f'(x)$.
    \item Να βρείτε το κρίσιμο σημείο συναρτήσει του $\lambda > 0$.
    \item Να βρείτε το $\lambda$ ώστε η $f$ να έχει μέγιστη τιμή $f(1) = -1$.
    \item Να επαληθεύσετε ότι πράγματι υπάρχει μέγιστο.
\end{alist}

\askhsh
Δίνεται η συνάρτηση $f(x) = x^3 + \alpha x^2 + \beta x - 2$.
\begin{alist}
    \item Να βρείτε την $f'(x)$.
    \item Να βρείτε τα $\alpha, \beta$ ώστε η $C_f$ να εφάπτεται στον άξονα $x'x$ στο $A(1, 0)$.
    \item Για τις τιμές αυτές, να εξετάσετε αν η $f$ έχει ακρότατο στο $x = 1$.
    \item Να βρείτε τα υπόλοιπα ακρότατα, αν υπάρχουν.
\end{alist}

\askhsh
Δίνεται η συνάρτηση $f(x) = x^3 + \alpha x^2 + 3x + 1$.
\begin{alist}
    \item Να βρείτε την $f'(x)$.
    \item Να βρείτε τη διακρίνουσα $\Delta$ της $f'(x) = 0$.
    \item Να βρείτε τις τιμές του $\alpha$ ώστε η $f$ να μην έχει ακρότατα.
    \item Να δείξετε ότι αν $|\alpha| < 3$, η $f$ είναι γνησίως αύξουσα στο $\mathbb{R}$.
\end{alist}

\askhsh
Δίνεται η συνάρτηση $f(x) = \alpha \sqrt{x} - x, \quad x > 0$.
\begin{alist}
    \item Να βρείτε την $f'(x)$.
    \item Να βρείτε το κρίσιμο σημείο συναρτήσει του $\alpha > 0$.
    \item Να βρείτε το $\alpha$ ώστε η $f$ να έχει μέγιστο στο $x_0 = 4$.
    \item Να υπολογίσετε τη μέγιστη τιμή.
\end{alist}

%=============================================
% ΟΜΑΔΑ Ε: ΑΠΟΔΕΙΞΕΙΣ ΑΝΙΣΟΤΗΤΩΝ
%=============================================
\paragraph{Ομάδα Ε: Αποδείξεις Ανισοτήτων με Ακρότατα}

\askhsh
Δίνεται η συνάρτηση $f(x) = e^x - x - 1$.
\begin{alist}
    \item Να βρείτε το πεδίο ορισμού της $f$.
    \item Να μελετήσετε τη συνάρτηση ως προς τη μονοτονία και τα ακρότατα.
    \item Να δείξετε ότι $e^x \ge x + 1$ για κάθε $x \in \mathbb{R}$.
    \item Να υπολογίσετε το $\displaystyle\lim_{x \to 0} \frac{e^x - x - 1}{x^2}$.
    \item Να εξετάσετε αν η $f$ είναι κυρτή ή κοίλη.
\end{alist}

\askhsh
Δίνεται η συνάρτηση $g(x) = x - \ln x - 1$, $x > 0$.
\begin{alist}
    \item Να μελετήσετε τη μονοτονία και τα ακρότατα της $g$.
    \item Να δείξετε ότι $\ln x \le x - 1$ για κάθε $x > 0$.
    \item Να βρείτε το $\displaystyle\lim_{x \to +\infty} \frac{\ln x}{x}$ χρησιμοποιώντας την ανισότητα.
    \item Να λύσετε την εξίσωση $\ln x = x - 1$.
    \item Να σχεδιάσετε τη γραφική παράσταση της $g$ σε κατάλληλο σύστημα αξόνων.
\end{alist}

\askhsh
Δίνεται η συνάρτηση $h(x) = x^2 - 2\ln x$, $x > 0$.
\begin{alist}
    \item Να βρείτε το πεδίο ορισμού και να υπολογίσετε τα $\displaystyle\lim_{x \to 0^+} h(x)$ και $\displaystyle\lim_{x \to +\infty} h(x)$.
    \item Να μελετήσετε τη μονοτονία και τα ακρότατα.
    \item Να δείξετε ότι $x^2 \ge 2\ln x + 1$ για κάθε $x > 0$.
    \item Να εξετάσετε αν η εξίσωση $x^2 = 2\ln x + 2$ έχει λύση στο $(0, +\infty)$.
\end{alist}

\askhsh
Δίνεται η συνάρτηση $f(x) = \sin x - x$.
\begin{alist}
    \item Να βρείτε το πεδίο ορισμού και να δείξετε ότι η $f$ είναι περιττή.
    \item Να μελετήσετε τη μονοτονία της $f$ στο $[0, +\infty)$.
    \item Να δείξετε ότι $\sin x < x$ για κάθε $x > 0$.
    \item Να υπολογίσετε το $\displaystyle\lim_{x \to 0} \frac{x - \sin x}{x^3}$.
    \item Χρησιμοποιώντας το ΘΜΤ, να δείξετε ότι $|\sin a - \sin b| \le |a - b|$ για κάθε $a, b \in \mathbb{R}$.
\end{alist}

\askhsh
Δίνεται η συνάρτηση $g(x) = x - 2\sqrt{x} + 1$, $x \ge 0$.
\begin{alist}
    \item Να δείξετε ότι $g(x) = (\sqrt{x} - 1)^2$.
    \item Να αποδείξετε ότι $2\sqrt{x} \le x + 1$ για κάθε $x \ge 0$.
    \item Για ποια τιμή του $x$ επέρχεται η ισότητα;
    \item Να βρείτε τις ασύμπτωτες της $C_g$.
    \item Να λύσετε την εξίσωση $\sqrt{x} = \frac{x+1}{2}$.
\end{alist}

\askhsh
Δίνεται η συνάρτηση $f(x) = e^x - 2x$.
\begin{alist}
    \item Να μελετήσετε τη μονοτονία και τα ακρότατα της $f$.
    \item Να υπολογίσετε τα $\displaystyle\lim_{x \to -\infty} f(x)$ και $\displaystyle\lim_{x \to +\infty} f(x)$.
    \item Να βρείτε το πλήθος των ριζών της εξίσωσης $e^x = 2x$.
    \item Να δείξετε ότι $e^x \ge 2x$ για κάθε $x \in \mathbb{R}$ και να βρείτε πότε ισχύει η ισότητα.
    \item Αν η εξίσωση έχει ρίζα $\rho$, να δείξετε ότι $\rho \in (0, 1)$.
\end{alist}



%=============================================
% ΟΜΑΔΑ ΣΤ: ΣΥΝΟΛΟ ΤΙΜΩΝ ΚΑΙ ΡΙΖΕΣ
%=============================================
\paragraph{Ομάδα ΣΤ: Σύνολο Τιμών και Πλήθος Ριζών}

\askhsh
Δίνεται η συνάρτηση $f(x) = x^2 + \frac{2}{x}, \quad x > 0$.
\begin{alist}
    \item Να υπολογίσετε την $f'(x)$.
    \item Να βρείτε τα κρίσιμα σημεία.
    \item Να μελετήσετε τη μονοτονία και τα ακρότατα.
    \item Να υπολογίσετε τα όρια στο $0^+$ και στο $+\infty$.
    \item Να βρείτε το σύνολο τιμών της $f$ στο $(0, +\infty)$.
\end{alist}

\askhsh
Δίνεται η συνάρτηση $f(x) = x^x, \quad x > 0$.
\begin{alist}
    \item Γράψτε την $f$ στη μορφή $f(x) = e^{x \ln x}$.
    \item Υπολογίστε την $f'(x)$.
    \item Βρείτε τα κρίσιμα σημεία.
    \item Βρείτε την ελάχιστη τιμή της $f$.
\end{alist}

\askhsh
Δίνεται η συνάρτηση $f(x) = x^3 - 3x + 5$.
\begin{alist}
    \item Να υπολογίσετε την $f'(x)$.
    \item Να βρείτε τα τοπικά ακρότατα.
    \item Να υπολογίσετε τις τιμές της $f$ στα σημεία των ακροτάτων.
    \item Να βρείτε το πλήθος των ριζών της εξίσωσης $f(x) = 0$.
\end{alist}

\askhsh
Δίνεται η συνάρτηση $f(x) = x^3 - 3x + 1$.
\begin{alist}
    \item Να μελετήσετε τη μονοτονία και τα ακρότατα.
    \item Να υπολογίσετε τις τιμές στα ακρότατα.
    \item Να βρείτε το πλήθος των ριζών της $f(x) = 0$.
    \item Να προσδιορίσετε τα διαστήματα στα οποία βρίσκονται οι ρίζες.
\end{alist}

\askhsh
Δίνεται η συνάρτηση $f(x) = \sin x + \cos x$ στο $[0, 2\pi]$.
\begin{alist}
    \item Να υπολογίσετε την $f'(x)$.
    \item Να βρείτε τα κρίσιμα σημεία στο $[0, 2\pi]$.
    \item Να βρείτε τη μέγιστη και την ελάχιστη τιμή της $f$.
    \item Να βρείτε το σύνολο τιμών της $f$ στο $[0, 2\pi]$.
\end{alist}

\askhsh
Να εξετάσετε αν οι παρακάτω εξισώσεις έχουν λύση:
\begin{alist}
    \item $e^x = x$
    \item $\ln x = x - 2$
    \item $x^3 - 3x + 4 = 0$
    \item $x^4 - 4x + 1 = 0$
\end{alist}

%=============================================
% ΟΜΑΔΑ Ζ: ΠΡΟΒΛΗΜΑΤΑ ΒΕΛΤΙΣΤΟΠΟΙΗΣΗΣ
%=============================================
\paragraph{Ομάδα Ζ: Προβλήματα Βελτιστοποίησης}

\askhsh
Από όλα τα ορθογώνια με περίμετρο $40$ cm:
\begin{alist}
    \item Αν $x, y$ οι διαστάσεις, να εκφράσετε το $y$ συναρτήσει του $x$.
    \item Να εκφράσετε το εμβαδόν $E$ ως συνάρτηση του $x$.
    \item Να βρείτε το πεδίο ορισμού της $E(x)$.
    \item Να βρείτε τις διαστάσεις που μεγιστοποιούν το εμβαδόν.
    \item Να υπολογίσετε το μέγιστο εμβαδόν.
\end{alist}

\askhsh
Να βρείτε δύο θετικούς αριθμούς με άθροισμα $10$ και μέγιστο γινόμενο.
\begin{alist}
    \item Αν $x, y$ οι αριθμοί, να εκφράσετε το $y$ συναρτήσει του $x$.
    \item Να εκφράσετε το γινόμενο $P$ ως συνάρτηση του $x$.
    \item Να βρείτε το πεδίο ορισμού.
    \item Να βρείτε τους δύο αριθμούς.
\end{alist}

\askhsh
Ένα ανοιχτό κουτί με τετράγωνη βάση έχει όγκο $32 \text{ m}^3$.
\begin{alist}
    \item Αν $x$ η πλευρά της βάσης και $h$ το ύψος, να εκφράσετε το $h$ συναρτήσει του $x$.
    \item Να εκφράσετε την επιφάνεια $E$ ως συνάρτηση του $x$.
    \item Να βρείτε τις διαστάσεις που ελαχιστοποιούν την επιφάνεια.
    \item Να υπολογίσετε την ελάχιστη επιφάνεια.
\end{alist}

\askhsh
Να βρείτε το σημείο της ευθείας $y = 2x + 3$ που απέχει ελάχιστη απόσταση από την αρχή $O(0,0)$.
\begin{alist}
    \item Να εκφράσετε την απόσταση $d$ ενός σημείου $(x, 2x+3)$ από το $O$.
    \item Να μελετήσετε τη συνάρτηση $D(x) = d^2(x)$ για ευκολία.
    \item Να βρείτε το $x$ που ελαχιστοποιεί την απόσταση.
    \item Να βρείτε το ζητούμενο σημείο.
\end{alist}

\askhsh
Δίνεται η συνάρτηση $f(x) = -x^2 + 4x - 3$.
\begin{alist}
    \item Να βρείτε την εξίσωση της εφαπτομένης στο σημείο $(x_0, f(x_0))$.
    \item Να βρείτε τα σημεία τομής της εφαπτομένης με τους άξονες.
    \item Να εκφράσετε το εμβαδόν του τριγώνου ως συνάρτηση του $x_0$.
    \item Να βρείτε την εφαπτομένη που σχηματίζει τρίγωνο ελάχιστου εμβαδού.
\end{alist}

\askhsh
Να βρείτε το σημείο της καμπύλης $y = \ln x$ που απέχει ελάχιστη απόσταση από την αρχή.
\begin{alist}
    \item Να εκφράσετε την απόσταση $d$ του σημείου $(x, \ln x)$ από το $O$.
    \item Να μελετήσετε τη συνάρτηση $D(x) = d^2(x)$.
    \item Να βρείτε (αριθμητικά ή αλγεβρικά) το ελάχιστο.
    \item Σχολιάστε τη μέθοδο.
\end{alist}

%=============================================
% ΟΜΑΔΑ Η: ΣΩΣΤΟ - ΛΑΘΟΣ
%=============================================
\paragraph{Ομάδα Η: Θεωρητικές Ερωτήσεις}

\askhsh
Να χαρακτηρίσετε τις παρακάτω προτάσεις ως Σωστές ή Λανθασμένες:
\begin{alist}
    \item Αν $f'(x_0) = 0$, τότε η $f$ έχει ακρότατο στο $x_0$.
    \item Αν η $f$ έχει τοπικό μέγιστο στο $x_0$, τότε $f'(x_0) = 0$.
    \item Αν η $f$ είναι γνησίως αύξουσα, τότε δεν έχει τοπικά ακρότατα.
    \item Τα ακρότατα μιας συνάρτησης βρίσκονται πάντα σε σημεία μηδενισμού της παραγώγου.
\end{alist}

\askhsh
Να χαρακτηρίσετε τις παρακάτω προτάσεις ως Σωστές ή Λανθασμένες:
\begin{alist}
    \item Μια συνεχής συνάρτηση σε κλειστό διάστημα έχει πάντα ολικό μέγιστο και ελάχιστο.
    \item Αν $f'(x) > 0$ για κάθε $x \in \mathbb{R}^*$, τότε η $f$ είναι γνησίως αύξουσα στο $\mathbb{R}$.
    \item Το ολικό μέγιστο μιας συνάρτησης είναι μοναδικό ως τιμή.
    \item Η θέση του ολικού μεγίστου είναι πάντα μοναδική.
\end{alist}

\askhsh
Δίνεται παραγωγίσιμη συνάρτηση $f: \mathbb{R} \to \mathbb{R}$ με $f'(x) > 0$ για κάθε $x \neq 0$.
\begin{alist}
    \item Είναι δυνατόν η $f$ να έχει ακρότατο στο $x = 0$; Αιτιολογήστε.
    \item Τι μπορείτε να πείτε για τη μονοτονία της $f$;
    \item Δώστε ένα παράδειγμα τέτοιας συνάρτησης.
    \item Ισχύει το ίδιο αν η $f$ δεν είναι συνεχής στο $0$;
\end{alist}

\askhsh
Να εξετάσετε αν οι παρακάτω συναρτήσεις έχουν ακρότατα:
\begin{alist}
    \item $f(x) = x^3$
    \item $f(x) = x^4$
    \item $f(x) = |x|$
    \item $f(x) = e^x$
\end{alist}

\end{document}
